\documentclass[11pt]{article}

\usepackage{nmrdoc}
\usepackage{siunitx}
\usepackage{bm}

\sisetup{per-mode=symbol}
\DeclareSIUnit{\angstrom}{\text{\AA}}

\begin{document}
\NmrTitle{The Dispersion Calculator}

\section{What the calculator computes}

The Dispersion calculator describes the short-range van der Waals geometry
around a nucleus, restricted in this implementation to nearby aromatic ring
vertices. London dispersion is the part of the intermolecular interaction that
comes from correlated electronic fluctuations. In a pair-potential formula it
is often represented by an attractive energy proportional to
\(-C_6/r^6\), where \(r\) is the separation and \(C_6\) carries the
species-dependent polarizability information.

The code does not compute that energy. It keeps the \(r^{-6}\) distance law
and the direction from a ring atom to the target nucleus, sets the dispersion
coefficient to one, and records the resulting geometry as a rank-2 tensor with
units \(\si{\angstrom^{-6}}\). ``Rank-2'' here means that the quantity has two
Cartesian indices: entry \(ab\) says how the geometry distinguishes an output
direction \(a\) from an input direction \(b\). The calculator also stores a
separate positive scalar \(r^{-6}\) contact strength, again with units
\(\si{\angstrom^{-6}}\).

This bears on NMR shielding because nearby polarizable electrons change the
electronic response of the observed nucleus to an applied magnetic field. The
calculator does not solve that electronic response. Its output is a geometric
descriptor that can correlate with a DFT shielding contribution after
calibration. Before calibration it is not a chemical shift and not a shielding
in ppm.

\section{Where shielding enters}

For a nucleus in an applied magnetic field \(\bm B_0\), the local field is
written
\begin{equation}
  \bm B_{\mathrm{nuc}} = (\bm I-\bm\sigma)\bm B_0 .
  \label{eq:shielding}
\end{equation}
Here \(\bm I\) is the \(3\times3\) identity matrix and \(\bm\sigma\) is the
nuclear shielding tensor. In physical NMR, \(\bm\sigma\) is dimensionless: it
is the linear-response coefficient connecting the applied field to the field
screened by the electrons. The matrix is needed because a molecule is not
spherically symmetric. A field applied along \(x\) can induce a secondary field
with \(x\), \(y\), and \(z\) components at the nucleus, and rotating the
molecule changes those components.

Equation~\eqref{eq:shielding} is a quantum-mechanical statement about
electrons. The external field perturbs the electronic state; the perturbed
electrons generate their own magnetic field at the nucleus; the part linear in
\(\bm B_0\) is recorded in \(\bm\sigma\). The isotropic chemical shift reported
in solution NMR is obtained from the orientationally averaged shielding after
reference subtraction and sign convention. That scalar observable is connected
to the tensor, but it does not contain all the tensor information.

The dispersion model used here isolates one geometric source of shielding
variation: how close aromatic ring atoms sit to the observed nucleus, and from
which directions they approach it. Aromatic rings are treated as sets of
vertices at their ring-atom coordinates. Each accepted vertex contributes a
short-range tensor whose magnitude falls as \(r^{-6}\). The result is not a
London dispersion energy and not an electronic-structure calculation. It is the
orientation-resolved geometry of nearby van der Waals contacts, left for
calibration to connect to the DFT shielding tensor.

\section{The tensor split used by the code}

The C++ type \texttt{SphericalTensor} decomposes any real \(3\times3\) tensor
\(\bm X\) into three pieces. The scalar piece is
\begin{equation}
  T_0 = \frac{1}{3}\operatorname{tr}\bm X ,
  \label{eq:t0}
\end{equation}
where \(\operatorname{tr}\bm X=X_{xx}+X_{yy}+X_{zz}\). This is the isotropic
part, the part that remains when all molecular orientations are averaged.

The antisymmetric part is stored as the three-component vector
\begin{equation}
  \bm T_1 =
  \frac{1}{2}
  \begin{pmatrix}
    X_{yz}-X_{zy}\\
    X_{zx}-X_{xz}\\
    X_{xy}-X_{yx}
  \end{pmatrix}.
  \label{eq:t1}
\end{equation}
This vector is a compact storage form for the part of the matrix that changes
sign under transpose. It is usually not measured directly in standard isotropic
solution NMR, but the full nine-number layout includes it.

The remaining part is the symmetric traceless matrix
\begin{equation}
  \bm S = \frac{\bm X+\bm X^{\mathsf T}}{2} - T_0\bm I,
  \qquad \operatorname{tr}\bm S=0 .
  \label{eq:symmetric-traceless}
\end{equation}
It has five independent entries. The code packs them as
\begin{equation}
  \bm T_2 =
  \begin{pmatrix}
    \sqrt{2}\,S_{xy}\\
    \sqrt{2}\,S_{yz}\\
    \sqrt{3/2}\,S_{zz}\\
    \sqrt{2}\,S_{xz}\\
    (S_{xx}-S_{yy})/\sqrt{2}
  \end{pmatrix}.
  \label{eq:t2}
\end{equation}
These factors are the normalization in \texttt{Types.cpp}. They preserve the
Frobenius length: \(\sum_m T_{2m}^2=\sum_{ab}S_{ab}^2\). The packed order is
\[
  [T_0,T_{1x},T_{1y},T_{1z},
    T_{2,-2},T_{2,-1},T_{2,0},T_{2,+1},T_{2,+2}] .
\]

For this calculator, the tensor contribution from each accepted vertex is
symmetric and traceless by construction. Its total therefore has
\(T_0=0\) and \(\bm T_1=\bm0\), up to floating-point roundoff, and its tensor
signal lives in the five \(T_2\) components. The separate scalar contact sum
is not the trace of that tensor. It is the stored isotropic \(r^{-6}\) contact
strength used as a separate feature.

\section{The method implemented}

Let target atom \(i\) have position \(\bm r_i\), in \(\si{\angstrom}\). Let an
aromatic ring \(R\) have vertex positions \(\bm v_j\), also in
\(\si{\angstrom}\), one for each ring atom. The code first queries ring centers
within \(\SI{15.0}{\angstrom}\), using the shared ring-current spatial cutoff
from \texttt{data/calculator\_params.toml}. This broad query only selects
candidate rings. The actual dispersion range is enforced per vertex by the
\(\SI{5.0}{\angstrom}\) cutoff below.

For each candidate ring, two ring-level exclusions are applied. The near-field
test rejects a target whose distance from the ring center is not greater than
the ring radius. In the code this filter uses the ring diameter as the source
extent and the parameter
\(\texttt{near\_field\_exclusion\_ratio}=0.5\). The topology test rejects the
whole ring if the target atom is a ring vertex or is covalently bonded to any
ring vertex. That removes through-bond cases from a through-space contact
model.

For an accepted ring and vertex \(j\),
\begin{equation}
  \bm d_{ij}=\bm r_i-\bm v_j,\qquad
  r_{ij}=|\bm d_{ij}|,\qquad
  \hat{\bm d}_{ij}=\frac{\bm d_{ij}}{r_{ij}} .
  \label{eq:vertex-geometry}
\end{equation}
The distance \(r_{ij}\) is in \(\si{\angstrom}\), while
\(\hat{\bm d}_{ij}\) is unitless. A vertex is skipped when
\(r_{ij}<\SI{0.1}{\angstrom}\), the singularity guard in the parameter file,
or when \(r_{ij}>\SI{5.0}{\angstrom}\), the dispersion vertex cutoff.

Between \(\SI{4.3}{\angstrom}\) and \(\SI{5.0}{\angstrom}\) the code multiplies
the contribution by a smooth switching function:
\begin{equation}
S(r)=
\begin{cases}
1, & r\le R_s,\\[3pt]
\dfrac{(R_c^2-r^2)^2(R_c^2+2r^2-3R_s^2)}
      {(R_c^2-R_s^2)^3}, & R_s<r<R_c,\\[9pt]
0, & r\ge R_c .
\end{cases}
\label{eq:switch}
\end{equation}
Here \(R_s=\SI{4.3}{\angstrom}\) is the switching onset and
\(R_c=\SI{5.0}{\angstrom}\) is the outer cutoff. Both values are read from
\texttt{data/calculator\_params.toml}. The function and its first derivative
match at both endpoints: the contribution is full strength before
\(R_s\), fades smoothly across the taper, and reaches zero at \(R_c\). If the
computed switch value is below \(10^{-15}\), the vertex is treated as zero.

The per-vertex tensor is
\begin{equation}
  K^{(ij)}_{ab}
  =
  S(r_{ij})
  \left(
    \frac{3\,d_{ij,a}d_{ij,b}}{r_{ij}^{8}}
    -\frac{\delta_{ab}}{r_{ij}^{6}}
  \right)
  =
  \frac{S(r_{ij})}{r_{ij}^{6}}
  \left(3\hat d_{ij,a}\hat d_{ij,b}-\delta_{ab}\right).
  \label{eq:kernel}
\end{equation}
The indices \(a,b\) run over \(x,y,z\). The Kronecker delta
\(\delta_{ab}\) is one when \(a=b\) and zero otherwise. Since
\(\bm d\bm d/r^8\) has units \(\si{\angstrom^{-6}}\), every entry of
\(\bm K^{(ij)}\) has units \(\si{\angstrom^{-6}}\). Its trace is zero:
the three diagonal entries sum to
\[
  \frac{S(r)}{r^6}\left(3|\hat{\bm d}|^2-3\right)=0 .
\]
The scalar contact stored beside it is
\begin{equation}
  q^{(ij)}=\frac{S(r_{ij})}{r_{ij}^{6}},
  \label{eq:scalar}
\end{equation}
also in \(\si{\angstrom^{-6}}\).

A useful exact picture is the eigendecomposition of
Eq.~\eqref{eq:kernel}. Each accepted vertex contributes a signed uniaxial
ellipsoid. Along the vertex-to-target direction \(\hat{\bm d}\), the tensor
has eigenvalue \(2S/r^6\). In either perpendicular direction it has eigenvalue
\(-S/r^6\). The signed sum of those three eigenvalues is zero, so the tensor
records anisotropy rather than volume. Moving the vertex farther away shrinks
all three axes as \(r^{-6}\), and the switching function scales all three axes
together. Summing a ring means adding these oriented, traceless ellipsoids
from its vertices.

For a ring \(R\), the code accumulates
\begin{equation}
  \bm K^{(iR)}=\sum_{j\in R_i} \bm K^{(ij)},\qquad
  q^{(iR)}=\sum_{j\in R_i} q^{(ij)},
  \label{eq:ring-sum}
\end{equation}
where \(R_i\) is the set of vertices of ring \(R\) that pass the distance,
cutoff, taper, and numerical-floor tests for target \(i\). If no vertex passes,
the ring contributes nothing. Otherwise the ring tensor is decomposed with
Eqs.~\eqref{eq:t0}--\eqref{eq:t2}; its \(T_2\) components and scalar contact
strength are added to the per-ring-type accumulators.

The atom total is then
\begin{equation}
  \bm K_i=\sum_{R\in\mathcal A_i}\bm K^{(iR)},
  \label{eq:atom-sum}
\end{equation}
where \(\mathcal A_i\) is the set of accepted rings with at least one accepted
vertex for atom \(i\). The code stores
\(\texttt{SphericalTensor::Decompose}(\bm K_i)\) in the field named
\texttt{disp\_shielding\_contribution}. The name is historical shorthand in the
data structure. The stored tensor is the uncalibrated
\(\si{\angstrom^{-6}}\) geometry from Eq.~\eqref{eq:atom-sum}.

The same per-vertex expression is used when the result is sampled at an
arbitrary point in space. That path skips points inside a ring-radius sphere,
points whose ring center is farther than \(\SI{15.0}{\angstrom}\), and
vertices that fail the same \(\SI{0.1}{\angstrom}\) singularity guard,
switching-floor test, or \(\SI{5.0}{\angstrom}\) vertex cutoff. It returns
the decomposed tensor field at the sample point and does not produce the
scalar contact output.

The main approximations are visible in these equations. The source sites are
aromatic ring vertices, not all atoms in the protein. Each vertex has unit
dispersion coefficient; element-specific \(C_6\) values and polarizabilities
are not multiplied into the tensor. The physical long-range \(r^{-6}\) tail is
truncated, with a taper over the last \(\SI{0.7}{\angstrom}\). The near-field
and through-bond exclusions define where the through-space vertex model is not
evaluated.

\section{Inputs and output}

This calculator consumes molecular geometry and topology: atom coordinates in
\(\si{\angstrom}\), aromatic ring membership, ring vertex coordinates, ring
centers and radii, and the covalent bond graph used for the through-bond
exclusion. Its declared runtime dependencies are \texttt{GeometryResult} and
\texttt{SpatialIndexResult}. It does not read MOPAC charges, AIMNet2 charges or
embeddings, ff14SB charges, APBS electrostatic fields, or DSSP hydrogen-bond
assignments; those are inputs to other calculators.

The calculator's own \texttt{WriteFeatures} method writes three arrays.

\texttt{disp\_shielding.npy} has shape \((N,9)\), one row per atom, in the
\texttt{SphericalTensor} order given above; for this tensor the \(T_2\)
entries carry the angular information.

\texttt{disp\_per\_type\_T0.npy} has shape \((N,8)\), with one scalar
\(q=\sum S/r^6\) contact sum per aromatic ring type. Despite the file suffix,
this scalar is not the trace of \(\bm K\).

\texttt{disp\_per\_type\_T2.npy} has shape \((N,40)\), with five \(T_2\)
components for each of the eight aromatic ring-type blocks. All three outputs
are geometric quantities in \(\si{\angstrom^{-6}}\).
Calibration against the DFT reference is what turns these descriptors into a
shielding contribution; before that step they should be read as geometry that
may correlate with the dispersion part of shielding.

\end{document}
